\section{Evaluation}

\subsection{Evaluation Setup}
We evaluate all methods under a strictly matched protocol.
Passenger demand is constructed from the NYC yellow taxi trip dataset (November 2025)\cite{nyc_yellow_taxi_dataset}.
The simulation window is fixed to 06:00--22:00 with 15-minute slots and receding horizon $H=64$.
Because real battery swapping station operation logs are not available in this study, station-side operations are simulated with parameterized station resources (battery inventory, charging slots, and swap service capacity).
All methods use the same demand input, seed policy, and three predefined resource settings.
Evaluation dates are fixed to 2025-11-01, 2025-11-05, and 2025-11-09.

We compare the proposed \textbf{COP} (Coordinated Optimization Policy) with three baselines:
\textbf{FDB} (Fleet Dispatch Baseline),
\textbf{HDB} (Heuristic Dispatch Baseline), and
\textbf{CCB} (Capacity-Constrained Charging Baseline).
FDB performs fleet-level dispatch by solving the network-level repositioning optimization each slot and sends low-energy taxis to swapping stations; station charging follows first-come-first-served.
HDB applies per-taxi in-zone heuristics: each taxi serves local demand if energy is sufficient, and low-energy taxis request local swapping only when local full-battery stock is available; station charging also follows first-come-first-served.
We additionally define \textbf{CCB} (Capacity-Constrained Charging Baseline), which keeps fleet-level dispatch but constrains the maximum number of simultaneously active chargers (e.g., 90\% of nominal chargers).

The primary performance metrics are served passengers and daily peak charging power demand.

\begin{table}[t]
\centering
\footnotesize
\setlength{\tabcolsep}{6pt}
\renewcommand{\arraystretch}{1.05}
\caption{Performance comparison (means over matched runs).}
\label{tab:primary_aggregate}
\begin{tabular}{lccc}
\hline
Method & Served Passengers & Peak Charging Demand (kW) & Peak Station Total Power (kW) \\
\hline
COP & $77613.67 \pm 6455.29$ & 31496.67 & 28530.00 \\
FDB & $77442.56 \pm 6369.63$ & 30840.00 & 29783.33 \\
HDB & $36284.33 \pm 3285.06$ & 22103.33 & 24836.67 \\
CCB & TBD & TBD & TBD \\
\hline
\end{tabular}
\end{table}

\subsection{Results}
\textbf{Performance Comparison:}
As shown in Table~\ref{tab:primary_aggregate}, COP improves served passengers by $+113.9\%$ over HDB and by $+0.22\%$ over FDB.
For daily peak charging power demand, COP is close to FDB and higher than HDB, indicating that service gains are not achieved through demand reduction alone.
CCB results are pending corrected rerun and will be filled in this table.

\textbf{Overhead Analysis:}
To quantify operational overhead, we report idle driving distance and battery handling frequency (number of swaps).

\begin{table}[t]
\centering
\footnotesize
\setlength{\tabcolsep}{5pt}
\renewcommand{\arraystretch}{1.05}
\caption{Overhead comparison (means over matched runs).}
\label{tab:overhead}
\begin{tabular}{lcc}
\hline
Method & Idle Driving Distance & Number of Swaps \\
\hline
COP & 38107.11 & 1946.11 \\
FDB & 37949.33 & 1980.67 \\
HDB & 0.00 & 781.22 \\
CCB & TBD & TBD \\
\hline
\end{tabular}
\end{table}

\textbf{Ablation Study:}
We separate taxi-side and station-side contributions under matched protocol settings.

\begin{table}[t]
\centering
\footnotesize
\setlength{\tabcolsep}{4pt}
\renewcommand{\arraystretch}{1.05}
\caption{Ablation contrasts (served passengers and peak demand).}
\label{tab:ablation_contrast}
\begin{tabular}{lrr}
\hline
Contrast & $\Delta$ Served & $\Delta$ Peak Demand (kW) \\
\hline
Variant A (COP vs FDB) & +171.1 & +656.7 \\
Variant B (Ideal dispatch with optimized vs FCFS charging) & +37.7 & -1060.0 \\
Dispatch effect under optimized charging & +41439.1 & +9383.3 \\
Dispatch effect under FCFS charging & +41158.2 & +8736.7 \\
\hline
\end{tabular}
\end{table}

\begin{table}[t]
\centering
\footnotesize
\setlength{\tabcolsep}{5pt}
\renewcommand{\arraystretch}{1.05}
\caption{COP solver runtime summary.}
\label{tab:runtime}
\begin{tabular}{lc}
\hline
Metric & Value \\
\hline
Avg reposition step runtime (mean $\pm$ std) & $10.20 \pm 0.52$ s \\
Reposition step runtime max (mean over runs) & 16.80 s \\
Worst observed per-run max & 21.79 s \\
\hline
\end{tabular}
\end{table}

CCB is treated as a supplementary capacity-constrained scenario and will not be merged into the primary paired-significance conclusion.

\subsection{Protocol-Aligned Time-Series Check}
Figure~\ref{fig:ts_panel} shows a protocol-aligned replay on 2025-11-05
under one representative resource setting, comparing COP and HDB.
The cumulative served difference is $+44947$ (COP minus HDB), consistent with aggregate results.
This replay is used for mechanism illustration, not as a replacement for paired statistical inference.

\begin{figure}[t]
  \centering
  \includegraphics[width=0.98\linewidth]{figures/evaluation/fig1_primary_served.pdf}
  \caption{Primary comparison on total\_served.}
  \label{fig:primary_served}
\end{figure}

\begin{figure}[t]
  \centering
  \begin{minipage}{0.32\linewidth}
    \centering
    \includegraphics[width=\linewidth]{figures/evaluation/fig3a_timeseries_served.pdf}\\[-1mm]
    \scriptsize (a) Served
  \end{minipage}\hfill
  \begin{minipage}{0.32\linewidth}
    \centering
    \includegraphics[width=\linewidth]{figures/evaluation/fig3b_timeseries_waiting.pdf}\\[-1mm]
    \scriptsize (b) Waiting
  \end{minipage}\hfill
  \begin{minipage}{0.32\linewidth}
    \centering
    \includegraphics[width=\linewidth]{figures/evaluation/fig3c_timeseries_station_power.pdf}\\[-1mm]
    \scriptsize (c) Station power
  \end{minipage}
  \caption{Protocol-aligned time-series check (COP vs HDB).}
  \label{fig:ts_panel}
\end{figure}
